
%\needspace{4\baselineskip}
\color{uniblue}{\subsection{Фиксация положения выключателей}}
\color{black}

\begin{enumerate}[label=\arabic{section}.\arabic{subsection}.\arabic{enumi}, labelsep=4pt, leftmargin=0pt, itemindent=57pt, itemsep=0pt, parsep=5pt] % если что можно крутить itemsep
%%%%%%%%%%%%%%%%%%%%%%%%%%%%%%%%%%%%%%%%%%%%%%%%%%%%%%%%%%%%%%%%%%%%%%  ОБЩИЕ СВЕДЕНИЯ %%%%%%%%%%%%%%%%%%%%%%%%%%%%%%%%%%%%%%%%%%%%%%%%%%

\item
Функциональный блок фиксации положения выключателей (ФПВ) предназначен для определения 
состояния выключателей на основании технологической информации:

\begin{itemize}
    \item состояния блок-контактов выключателя;
    \item положения разъединителей в цепи выключателя;
    \item информации о выводе в ремонт выключателя.
\end{itemize}

\item
Контроль положения разъединителей при формировании сигналов включенного или отключенного 
положения выключателей (например, \textit{«В1 включен»} и \textit{«В2 отключен»} для выключателя В1) вводится 
программным ключом \textbf{«Контр. Р В1»}. Также имеется возможность выбрать количество разъединителей, 
участвующих в формировании сигналов включенного или отключенного положения выключателей, для 
выключателя В1 она задается программным переключателем \textbf{«Схема Р В1»}.

\item
Также функция ФПВ пофазно контролирует состояние блок-контактов выключателя и в случае 
расхождения полюсов формирует сигнал \textit{«Внутр. пуск. ЗНФ В1»}.

\end{enumerate}